\section{Introduction}

Algebraic branching programs (ABPs) interpolate between formulas and general arithmetic
circuits.  A width-one ABP is simply a directed path, and hence computes a product of
affine linear forms.  Applying a Boolean hypercube summation to such products gives the
nondeterministic class \(\VNPone\).  Bringmann, Ikenmeyer and Zuiddam
\cite{BringmannIkenmeyerZuiddam2018,BringmannIkenmeyerZuiddam2017} proved
\[
  \VNPone(\F)=\VNP(\F)
  \quad\text{when }\charac(\F)\neq 2,
\]
and sharpened their construction so that each affine factor uses at most two variables.
Their argument uses identities containing division by two.  They showed separately that
\(\VNPone(\F_2)\subsetneq\VNP(\F_2)\), and explicitly left the remaining
characteristic-two fields open.

This paper supplies the missing characteristic-two construction whenever the field has an
element other than zero and one.  The exceptional field is therefore exactly \(\F_2\), in
agreement with the earlier separation.

\subsection{Main finite theorem}

We count the size of a binary arithmetic formula by the number of leaves, addition gates,
and multiplication gates.  Leaves are variables or field constants.

\begin{theorem}[Support-two characteristic-two compiler]\label{thm:finite-main}
Let \(\F\) be a field of characteristic two, let
\(\tau\in\F\setminus\{0,1\}\), and let \(\Phi\) be a binary division-free arithmetic
formula of size \(s\).  There exist integers \(q,M\ge 0\) and affine linear forms
\[
  L_1,\ldots,L_M\in
  \F[\mathbf{x},b_1,\ldots,b_q]
\]
such that each \(L_j\) involves at most two variables and
\[
  \Phi(\mathbf{x})
  =\sum_{\mathbf{b}\in\B^q}\prod_{j=1}^{M}L_j(\mathbf{x},\mathbf{b}).
\]
Moreover,
\[
  q\le 14s,
  \qquad
  M\le 34s.
\]
\end{theorem}

The formal artifact proves a stronger syntactic statement for its generated factors: each
factor contains at most two variable terms in its explicit affine-form term list.  This
implies the usual semantic support bound appearing in \cref{thm:finite-main}.

\subsection{Class consequences}

Write \(\VNPtwo(\F)\) for polynomial families admitting polynomial-dimensional Boolean
hypercube sums of polynomially many affine forms, each involving at most two variables.
Write \(\Ncl(\mathcal C)\) for the nondeterministic closure of a class \(\mathcal C\).

\begin{theorem}[Internal class equality]\label{thm:class-main}
Let \(\F\) have characteristic two and suppose \(\F\neq\F_2\).  Then
\[
   \Ncl(\VPe(\F))=\VNPtwo(\F).
\]
The statement covers finite and infinite fields.
\end{theorem}

A product of affine forms is a width-one ABP: place the factors consecutively on a directed
path.  Thus the support-two class maps directly into the class denoted
\(\VNPwplus\) by Bringmann--Ikenmeyer--Zuiddam.  Conversely, constant-width ABPs have
polynomial-size formulas, and Valiant's characterization gives
\(\Ncl(\VPe)=\VNP\) \cite{Valiant1980,BringmannIkenmeyerZuiddam2018}.

\begin{corollary}[Complete field classification]\label{cor:classification}
For every field \(\F\),
\[
  \VNPwplus(\F)=\VNP(\F)
  \quad\Longleftrightarrow\quad |\F|>2.
\]
For unrestricted affine labels,
\[
  \VNPone(\F)=\VNP(\F)
  \quad\Longleftrightarrow\quad \F\not\cong\F_2.
\]
\end{corollary}

For characteristic different from two, the first equality is the support-two theorem of
Bringmann--Ikenmeyer--Zuiddam.  For characteristic two and \(\F\neq\F_2\), it follows from
\cref{thm:class-main}.  Over \(\F_2\), even unrestricted width one is strictly weaker than
\(\VNP\) by their Proposition~9.1.

\subsection{Proof architecture}

The proof is deliberately modular.

\begin{enumerate}[label=\textbf{Step \arabic*.},leftmargin=2.2cm]
\item Convert a formula into a ranked, labelled acyclic graph whose path polynomial is the
formula.  The graph has maximum total degree three and labels involving at most one
original variable.
\item Introduce one Boolean selector variable for each edge.  A product of affine flow
factors and quadratic pair exclusions is the indicator of a selected source--sink path.
Multiplying by conditional edge-label factors and summing over all selectors recovers the
ABP path polynomial.
\item Replace every quadratic factor \(1+xy\) by a one-bit sum of three support-two affine
factors.
\item Replace each ternary internal flow factor by a one-bit sum of four support-two affine
factors.  The new identity carries an explicit cubic error.
\item Fix the original edge-selector assignment.  At every ternary vertex, two incident
edges share a head or a tail, so an already-present pair exclusion kills the cubic error
pointwise on the Boolean cube.  Only after this cancellation do we sum over the original
selectors.
\item Count the fresh auxiliaries and factors.  A sharp structural identity shows that the
number of ordered pair exclusions is four times the number of addition gates.
\item Flatten an outer \(\VNP\)-type witness cube with the compiler's fresh cube to obtain
the class-level inclusion.  The reverse inclusion follows because a product of
polynomially many affine forms has a polynomial-size formula.
\end{enumerate}

The distinction in Step~5 is not cosmetic.  The polynomial
\(abc(1+ab)\) is nonzero in \(\F[a,b,c]\); it vanishes only after \(a,b,c\) are specialized
to Boolean values.  The formal development contains an explicit negative theorem
preventing this pointwise argument from being silently promoted to a polynomial identity.

\subsection{Scope}

The result is a representation theorem.  The hypercube may have linearly many witness
bits, and no faster evaluator for the resulting sum is provided.  Consequently, nothing
here separates \(\VP\) from \(\VNP\), resolves a Boolean complexity-class problem, or
implies a subexponential algorithm.

The characteristic-two border-width-two problem studied by Dutta et al.
\cite{DuttaEtAl2024} is adjacent but different.  Their model is deterministic and
border-approximate; the present construction is exact and nondeterministic.  No
Frobenius-root extraction or border simulation is claimed here.

\section{Preliminaries}\label{sec:prelim}

\subsection{Polynomials, formulas, and families}

Throughout, \(\F\) is a field.  A \emph{binary arithmetic formula} is a rooted binary tree
whose internal nodes are labelled by \(+\) or \(\times\), and whose leaves are variables
or elements of \(\F\).  Its size is the number of nodes.  If \(l,a,u\) denote the numbers
of leaves, addition gates, and multiplication gates, respectively, then
\[
  s=l+a+u.
\]
A family \(f=(f_n)_{n\ge 0}\) is a \emph{p-family} if its number of variables and total
degree are polynomially bounded.  We use standard definitions of \(\VPe\) and \(\VNP\);
see, for example, \cite{Burgisser2000,MalodPortier2008}.  The class \(\VPe\) consists of
p-families with polynomially bounded formula size.

\subsection{Algebraic branching programs}

An algebraic branching program is a finite directed acyclic graph \(G=(V,E)\) with a
source \(s\), a sink \(t\), and an affine label \(\ell_e\) on every edge.  Its value is
\[
  \val(G)=\sum_{P:s\leadsto t}\prod_{e\in P}\ell_e.
\]
A width-one ABP is a directed path and hence computes a product of affine forms.  We use
\(\VPone\) for polynomial-size width-one ABPs with unrestricted affine labels, and
\(\VPone^{\mathrm{w}+}\) for the version in which every edge label involves at most two
variables.  Following Bringmann--Ikenmeyer--Zuiddam, set
\[
  \VNPone=\Ncl(\VPone),
  \qquad
  \VNPwplus=\Ncl(\VPone^{\mathrm{w}+}).
\]

\subsection{Nondeterministic closure}

\begin{definition}[Nondeterministic closure]\label{def:nclosure}
For a class \(\mathcal C\) of polynomial families, \(\Ncl(\mathcal C)\) consists of all
families \(f=(f_n)\) for which there exist a family \(g=(g_n)\in\mathcal C\) and
polynomially bounded functions \(p,q\) such that
\[
  f_n(\mathbf{x})
   =\sum_{\mathbf{b}\in\B^{p(n)}}g_{q(n)}(\mathbf{b},\mathbf{x}).
\]
\end{definition}

This is the closure operator introduced as Definition~2.2 in
\cite{BringmannIkenmeyerZuiddam2018}.  We will use its monotonicity:
\[
  \mathcal C\subseteq\mathcal D
  \quad\Longrightarrow\quad
  \Ncl(\mathcal C)\subseteq\Ncl(\mathcal D).
\]

\subsection{Affine-product hypercube representations}

\begin{definition}[Support-two affine-hypercube representation]\label{def:rep}
Let \(f\in\F[\mathbf{x}]\).  A support-two affine-hypercube representation of \(f\) is an
identity
\[
  f(\mathbf{x})
  =\sum_{\mathbf{b}\in\B^q}\prod_{j=1}^{M}L_j(\mathbf{x},\mathbf{b}),
\]
where every \(L_j\) is affine and \(|\supp(L_j)|\le 2\).  Constants are allowed, and the
support counts both original and hypercube variables.
\end{definition}

For families, \(q\) and \(M\) must be polynomially bounded.  The resulting class is denoted
\(\VNPtwo\).  Since a product of \(M\) affine factors is computed by a width-one path of
length \(M\), one has the immediate bridge
\[
  \VNPtwo\subseteq\VNPwplus.
\]

\section{Characteristic-two gadgets}\label{sec:gadgets}

Assume throughout this section that \(\charac(\F)=2\).  Fix
\(\tau\in\F\setminus\{0,1\}\) and put
\[
  \kappa=\tau(\tau+1).
\]

\begin{lemma}[Denominator legality]\label{lem:kappa}
The element \(\kappa\) is nonzero.
\end{lemma}

\begin{proof}
Both \(\tau\) and \(\tau+1\) are nonzero.  Indeed, \(\tau+1=0\) would imply
\(\tau=-1=1\) in characteristic two.  Since \(\F\) is a field, their product is nonzero.
\end{proof}

\subsection{A one-bit quadratic gadget}

\begin{lemma}[Quadratic identity]\label{lem:quadratic}
For arbitrary elements \(x,y\) of any commutative \(\F\)-algebra,
\[
\sum_{h\in\B}
  \kappa^{-1}(\tau+1+h)(x+\tau+h)(\kappa y+\tau+h)
 =1+xy.
\]
Only \(h\) is Boolean; no Boolean relation is imposed on \(x\) or \(y\).
\end{lemma}

\begin{proof}
First prove the denominator-free identity.  The left-hand side before multiplication by
\(\kappa^{-1}\) is
\[
 S=(\tau+1)(x+\tau)(\kappa y+\tau)
   +\tau(x+\tau+1)(\kappa y+\tau+1).
\]
The coefficient of \(\kappa y\) is
\[
 (\tau+1)(x+\tau)+\tau(x+\tau+1)
 =x\bigl((\tau+1)+\tau\bigr)
   +2\tau(\tau+1)=x.
\]
The remaining term is
\[
 \tau(\tau+1)\bigl((x+\tau)+(x+\tau+1)\bigr)=\kappa.
\]
Thus \(S=\kappa xy+\kappa=\kappa(1+xy)\).  Divide by \(\kappa\), using
\cref{lem:kappa}.
\end{proof}

The three factors in each summand involve, respectively, only \(h\), the pair \((x,h)\),
and the pair \((y,h)\), whenever \(x\) and \(y\) are affine expressions supported on one
variable each.  A distinct fresh \(h\) will be used for every occurrence of the gadget.

\subsection{A ternary gadget with a controlled error}

Define the two-point functional
\[
 \Lambda_\tau[p]
   =(\tau+1)p(\tau)+\tau p(\tau+1).
\]
A direct calculation in characteristic two gives
\[
 \Lambda_\tau[1]=1,
 \qquad
 \Lambda_\tau[T]=0,
 \qquad
 \Lambda_\tau[T^2]=\kappa,
 \qquad
 \Lambda_\tau[T^3]=\kappa.
\]

\begin{lemma}[Ternary identity]\label{lem:ternary}
For arbitrary elements \(a,b,c\) of any commutative \(\F\)-algebra,
\begin{align*}
&\sum_{h\in\B}
 \kappa^{-1}(\tau+1+h)(a+\tau+h)(b+\tau+h)(c+\tau+h)\\
&\hspace{35mm}=1+a+b+c+\kappa^{-1}abc.
\end{align*}
\end{lemma}

\begin{proof}
Apply \(\Lambda_\tau\) to
\[
 p(T)=(a+T)(b+T)(c+T)
 =abc+(ab+ac+bc)T+(a+b+c)T^2+T^3.
\]
The four moments displayed above yield
\[
 \Lambda_\tau[p]=abc+\kappa(a+b+c)+\kappa
 =\kappa(1+a+b+c)+abc.
\]
Division by \(\kappa\) proves the claim.
\end{proof}

The right-hand side contains the genuine cubic error \(\kappa^{-1}abc\).  It must not be
discarded as a formal polynomial.

\begin{lemma}[Boolean error killer]\label{lem:killer}
For \(u,v,w\in\B\), viewed inside any characteristic-two field,
\[
  uvw(1+uv)=0.
\]
\end{lemma}

\begin{proof}
If \(uv=0\), the product vanishes.  If \(uv=1\), then \(1+uv=1+1=0\).
\end{proof}

\begin{remark}[Formal versus pointwise cancellation]\label{rem:firewall}
The polynomial
\[
 abc(1+ab)=abc+a^2b^2c
\]
is nonzero in \(\F[a,b,c]\).  The cancellation in \cref{lem:killer} is valid only after
specializing \(a,b,c\) to Boolean values.  This distinction determines the proof order in
\cref{sec:compiler} and is separately machine-checked in the formal artifact.
\end{remark}
